% Generated draft: requires mathematical and source review.
\documentclass[11pt]{article}
\usepackage[margin=25mm]{geometry}
\usepackage{fontspec}
\setmainfont{Latin Modern Roman}
\usepackage{amsmath,amssymb,mathtools,unicode-math}
\setmathfont{Latin Modern Math}
\usepackage{enumitem,needspace,xurl,hyperref}
\setlength{\parindent}{0pt}
\setlength{\parskip}{5pt}
\setlength{\emergencystretch}{3em}
\newcommand{\N}{\mathbb{N}}
\newcommand{\Z}{\mathbb{Z}}
\newcommand{\Q}{\mathbb{Q}}
\newcommand{\R}{\mathbb{R}}
\newcommand{\C}{\mathbb{C}}
\newcommand{\F}{\mathbb{F}}
\newcommand{\A}{\mathbb{A}}
\newcommand{\PP}{\mathbb{P}}
\newcommand{\E}{\mathbb{E}}
\newcommand{\eps}{\varepsilon}
\newcommand{\dd}{\,\mathrm{d}}
\begin{document}
\section*{UnsolvedMath ID: 30001236}
\textbf{Source problem code:} catalogue entry
\subsection{Definitions and mathematical statement}
For a standard graded algebra \(R\) over a field \(k\), with \(R_0=k\), and a finitely generated graded \(R\)-module \(M\), the usual graded Betti numbers are
\[
 b^R_{i,j}(M)=\dim_k\operatorname{Tor}^R_i(M,k)_j,
 \qquad i\in\Z_{\geq 0},\quad j\in\Z.
\]
This definition explains the conventional notation; the supplied record does not explicitly specify the field, grading, or module conventions.

The source asks the following conditional question. With \(S,F,P,I,L\) satisfying the hypotheses referred to there as those of the Lex-Plus-Powers Conjecture, does one have
\[
 b_{i,j}^{S/P}(L)\geq b_{i,j}^{S/F}(I)
 \qquad\text{for all }i,j?
\]
Under the graded Betti-number convention above, the index range is \(i\in\Z_{\geq0}\) and \(j\in\Z\).

The referenced hypotheses are not reproduced in the supplied record. In particular, it does not define \(S,F,P,I,L\), specify the regular-sequence or degree conditions, state the relation between the Hilbert functions, or identify precisely the modules denoted by \(I\) and \(L\) over the respective quotient rings. Consequently, this is a transcription of the intended question with an unresolved contextual dependency, rather than a fully specified standalone conjecture. The arguments \(L\) and \(I\) in the displayed inequality are retained exactly; replacing them with quotient modules would require source clarification.
\subsection{Short English statement}
The question asks whether the lex-plus-powers object has at least as many graded Betti numbers over \(S/P\) as the corresponding object over \(S/F\), in every homological and internal degree. The supplied record omits the hypotheses defining these objects, so the precise scope and module interpretation require review.
\subsection{Sources}
\begin{itemize}
\item Supplied local UnsolvedMath registry record, permanent ID \(30001236\), title ``Relative Betti Inequalities for Lex-Plus-Powers Ideals,'' source problem code \texttt{catalogue entry}, in the Oberwolfach Reports Open Problems collection. The record gives the proposed year as 2009 and supplies no individual proposer.
\item Primary source supplied by the registry: \emph{Kommutative Algebra} (2009), Oberwolfach Reports, report 22, pp.\ 1157--1224, DOI 10.4171/owr/2009/22. \url{https://doi.org/10.4171/owr/2009/22}. The supplied background describes these reports as records of mathematical workshops and their open-problem discussions.
\item Supplied literature source: ``Caviglia-Sammartano on lex-plus-powers.'' \url{https://arxiv.org/abs/1802.03035}.
\item Supplied literature source: ``A study of the lex-plus-powers conjecture.'' \url{https://doi.org/10.1016/S0022-4049(03)00130-0}.
\item Literature assessment dated exactly \texttt{2026-08-21}: the supplied curation labels the problem open, describes it as the general Betti-number inequality in the lex-plus-powers setting, states that the report leaves it open outside the pure-power case, and says that later work treats only broad special families. This is supplied dated curation, not an independently verified literature assessment; its status and cited claims require independent verification.
\item The registry labels the statement \texttt{corrected\_verified} and reports that its mathematical content was checked against the official Oberwolfach report, with notation, context, and question form normalized without changing the source-level claim. This is supplied verification metadata, not a verification performed for this draft.
\end{itemize}
\end{document}
